\documentclass[11pt,a4paper]{article}
\usepackage[utf8]{inputenc}
\usepackage[T1]{fontenc}
\usepackage[french]{babel}
\usepackage[margin=2cm]{geometry}
\usepackage{enumitem}
\usepackage{booktabs}
\usepackage{array}
\usepackage{xcolor}
\usepackage{hyperref}
\usepackage{fancyhdr}
\usepackage{titlesec}

\definecolor{bleu}{HTML}{0b3d91}
\definecolor{orange}{HTML}{fd7e14}

\hypersetup{colorlinks=true,linkcolor=bleu,urlcolor=bleu}

\titleformat{\section}{\Large\bfseries\color{bleu}}{}{0em}{}[\vspace{-0.5em}\color{bleu}\rule{\textwidth}{0.5pt}]
\titleformat{\subsection}{\large\bfseries\color{bleu}}{}{0em}{}

\pagestyle{fancy}
\fancyhf{}
\fancyhead[L]{\small\color{gray}Rapport technique -- TP Examen 1}
\fancyhead[R]{\small\color{gray}Groupe 25}
\fancyfoot[C]{\small\thepage}
\renewcommand{\headrulewidth}{0.4pt}

\begin{document}

% ============ PAGE DE TITRE ============
\begin{titlepage}
\centering
\vspace*{3cm}
{\Huge\bfseries\color{bleu} Rapport Technique\par}
\vspace{0.5cm}
{\LARGE TP Examen 1\par}
\vspace{1cm}
{\Large\itshape Gestion d'une ligne de Commande\par}
\vspace{2cm}
{\large
\begin{tabular}{rl}
\textbf{Groupe :} & 25 \\
\textbf{Membres :} & Ngomde Cadet Kamdaou \\
                    & Bediang Leandrine Orline \\
\textbf{Formation :} & Master DI \\
\textbf{Institution :} & Universit\'e de Moundou \\
\textbf{Date :} & Octobre 2026 \\
\end{tabular}
\par}
\vfill
{\small Programmation Web -- Ann\'ee acad\'emique 2025--2026}
\end{titlepage}

\tableofcontents
\newpage

% ============ CONTENU ============

\section{Pr\'esentation du projet}

Ce projet impl\'emente une application web de gestion de commandes en PHP/MySQL.
Elle permet la cr\'eation, consultation, modification et suppression de produits,
clients et commandes, avec un tableau de bord statistique pour l'analyse de
performance des lignes de commande.

L'application r\'epond aux cinq t\^aches du sujet d'examen :
\begin{enumerate}
  \item Formulaire d'enregistrement des produits d'une commande
  \item Affichage des listes (produits, commandes, clients)
  \item Modification d'un produit
  \item Suppression d'un produit dans une commande
  \item Tableau de bord avec statistiques sur les lignes de commande
\end{enumerate}


\section{Architecture de la base de donn\'ees}

Le sch\'ema relationnel suit fid\`element le mod\`ele impos\'e par le sujet :

\begin{table}[h]
\centering
\begin{tabular}{llll}
\toprule
\textbf{Table} & \textbf{Cl\'e primaire} & \textbf{Cl\'es \'etrang\`eres} & \textbf{R\^ole} \\
\midrule
Departement     & idDep (AUTO\_INCREMENT)        & --                        & R\'ef\'erentiel d\'epartements \\
Region          & idRegion (AUTO\_INCREMENT)      & idDep                     & R\'ef\'erentiel r\'egions \\
Client          & idCli (AUTO\_INCREMENT)         & idRegion                  & Clients (photo, email unique) \\
Produit         & idProd (AUTO\_INCREMENT)        & --                        & Catalogue produits \\
Commande        & idCmd (AUTO\_INCREMENT)         & idCli                     & En-t\^ete de commande \\
ligneCommande   & (idCmd, idProd)                 & idCmd, idProd             & D\'etail : quantit\'e par produit \\
\bottomrule
\end{tabular}
\caption{Structure des tables}
\end{table}

\subsection{Choix techniques}

\begin{itemize}
  \item \textbf{ENGINE=InnoDB} : support des cl\'es \'etrang\`eres et des transactions ACID.
  \item \textbf{ON DELETE CASCADE} : la suppression d'une commande supprime automatiquement ses lignes, \'evitant les enregistrements orphelins.
  \item \textbf{UNIQUE KEY sur email} : l'unicit\'e des clients est garantie au niveau base de donn\'ees.
  \item \textbf{Charset utf8mb4} : support complet des caract\`eres Unicode (accents, caract\`eres sp\'eciaux).
  \item \textbf{Donn\'ees de test} : le fichier \texttt{schema.sql} ins\`ere 5 d\'epartements, 5 r\'egions, 5 clients, 8 produits, 5 commandes et 12 lignes de commande pour faciliter les tests.
\end{itemize}


\section{Architecture applicative}

\subsection{Organisation des fichiers}

\begin{table}[h]
\centering
\begin{tabular}{>{\ttfamily}l l}
\toprule
\textbf{Fichier} & \textbf{R\^ole} \\
\midrule
config.php              & Connexion PDO, fonctions utilitaires \\
includes/navbar.php     & Barre de navigation partag\'ee \\
index.php               & Tableau de bord (KPIs + graphiques) \\
produits.php            & Liste des produits avec recherche \\
produit\_form.php       & Formulaire cr\'eation/modification produit \\
produit\_save.php       & Traitement POST produit \\
produit\_delete.php     & Suppression produit \\
clients.php             & Liste des clients avec recherche \\
client\_form.php        & Formulaire cr\'eation/modification client \\
client\_save.php        & Traitement POST client \\
commandes.php           & Liste des commandes \\
commande\_form.php      & Formulaire nouvelle commande (multi-lignes) \\
commande\_save.php      & Traitement POST commande (transactionnel) \\
commande\_detail.php    & D\'etail d'une commande \\
ligne\_delete.php       & Suppression d'une ligne de commande \\
uploads/                & Stockage des photos \\
\bottomrule
\end{tabular}
\caption{Fichiers du projet}
\end{table}

\subsection{Pattern architectural}

L'application suit un pattern \textbf{page-controller} : chaque URL correspond \`a un fichier PHP qui g\`ere la logique m\'etier et le rendu. Les fichiers \texttt{*\_save.php} et \texttt{*\_delete.php} traitent exclusivement les requ\^etes POST et redirigent via \texttt{Location}, appliquant le pattern \textbf{POST/Redirect/GET (PRG)} pour \'eviter les soumissions doubles lors d'un rafra\^ichissement de page.


\section{S\'ecurit\'e}

La s\'ecurit\'e repose sur un principe fondamental : \textit{ne jamais faire confiance aux donn\'ees du client}.

\begin{itemize}
  \item \textbf{Requ\^etes pr\'epar\'ees PDO} : toutes les requ\^etes SQL utilisent des param\`etres li\'es (\texttt{?} ou \texttt{:placeholder}), \'eliminant le risque d'injection SQL.
  \item \textbf{Protection XSS} : la fonction \texttt{h()} applique \texttt{htmlspecialchars()} avec \texttt{ENT\_QUOTES} sur toute donn\'ee affich\'ee dans le HTML.
  \item \textbf{Validation serveur} : chaque fichier de traitement revalide les donn\'ees (champs obligatoires, format email, types num\'eriques) ind\'ependamment de la validation JavaScript.
  \item \textbf{M\'ethode HTTP} : les op\'erations de modification utilisent exclusivement POST ; les fichiers de traitement rejettent les requ\^etes GET avec un code 405.
  \item \textbf{Upload s\'ecuris\'e} : v\'erification de l'extension, renommage avec \texttt{uniqid()} pour \'eviter l'\'ecrasement et les attaques par nom de fichier.
\end{itemize}


\section{Interface utilisateur}

\subsection{Technologies front-end}

\begin{itemize}
  \item \textbf{Bootstrap 5.3} : framework CSS responsive, garantissant un affichage correct sur mobile et desktop via le syst\`eme de grille et les composants adaptatifs.
  \item \textbf{Bootstrap Icons} : ic\^ones vectorielles coh\'erentes dans toute l'interface.
  \item \textbf{Chart.js 4} : biblioth\`eque de graphiques l\'eg\`ere utilis\'ee pour le tableau de bord.
\end{itemize}

\subsection{Choix d'ergonomie}

\begin{itemize}
  \item \textbf{Messages flash} : notifications temporaires stock\'ees en session PHP, affich\'ees une seule fois puis effac\'ees. Cela \'evite les messages persistants apr\`es un rafra\^ichissement.
  \item \textbf{Calcul en temps r\'eel} : dans le formulaire de commande, les sous-totaux et le total se recalculent instantan\'ement via JavaScript \`a chaque modification de produit ou quantit\'e.
  \item \textbf{Badges de stock} : code couleur vert ($> 10$), orange ($\leq 10$), rouge ($= 0$) pour rep\'erer rapidement les ruptures.
  \item \textbf{Confirmation de suppression} : toute action destructrice demande une confirmation utilisateur via \texttt{confirm()}.
  \item \textbf{Navigation unifi\'ee} : barre de navigation partag\'ee avec d\'etection automatique de la page active.
\end{itemize}


\section{Tableau de bord -- Analyse de performance}

Le tableau de bord (page d'accueil) r\'epond \`a la t\^ache 5 du sujet en proposant une analyse multi-axes :

\begin{enumerate}
  \item \textbf{4 cartes KPI} : nombre de produits, clients, commandes, et chiffre d'affaires total. Offrent une vue synth\'etique imm\'ediate.
  \item \textbf{Graphique CA/mois} : diagramme \`a barres combin\'e avec une courbe du nombre de commandes (axe double). Permet de corr\'eler le volume de commandes avec le revenu g\'en\'er\'e et d'identifier les tendances.
  \item \textbf{Top 5 produits} : diagramme en anneau (doughnut) montrant les produits les plus command\'es en quantit\'e. Utile pour orienter les d\'ecisions d'approvisionnement.
  \item \textbf{Commandes r\'ecentes} : tableau des 5 derni\`eres commandes avec nombre de lignes et montant.
  \item \textbf{Alertes stock faible} : produits dont le stock est $\leq 10$, badge rouge pour stock \`a z\'ero. Permet d'anticiper les ruptures.
  \item \textbf{CA par client} : classement des clients par chiffre d'affaires g\'en\'er\'e avec nombre de commandes, permettant d'identifier les meilleurs clients.
\end{enumerate}

Cette approche permet une analyse de performance sur quatre axes : \textit{temporel} (tendance mensuelle), \textit{produit} (popularit\'e), \textit{client} (valeur), et \textit{logistique} (gestion des stocks).


\section{Gestion des transactions}

La cr\'eation d'une commande implique deux op\'erations :
\begin{enumerate}
  \item Insertion de l'en-t\^ete dans la table \texttt{Commande}
  \item Insertion de chaque ligne dans \texttt{ligneCommande}
\end{enumerate}

Ces op\'erations sont encapsul\'ees dans une \textbf{transaction} (\texttt{beginTransaction} / \texttt{commit} / \texttt{rollBack}), garantissant l'\textbf{atomicit\'e} : soit la commande compl\`ete est enregistr\'ee, soit rien n'est modifi\'e en cas d'erreur. Cela \'evite les commandes sans lignes ou les lignes orphelines.


\section{Instructions de d\'eploiement}

\begin{enumerate}
  \item Importer \texttt{schema.sql} dans phpMyAdmin avec les identifiants fournis (Login : \texttt{MasTMound2}, Mot de passe : \texttt{MoundMasT2@2026}).
  \item Placer le dossier \texttt{tpexam1} dans le r\'epertoire \texttt{htdocs} du serveur web (XAMPP, WAMP, etc.).
  \item V\'erifier que le dossier \texttt{uploads/} est accessible en \'ecriture par le serveur web.
  \item Ouvrir \texttt{http://localhost/tpexam1/} dans le navigateur.
\end{enumerate}

\vfill
\begin{center}
\small\color{gray}
Groupe 25 -- TP Examen Programmation Web -- Master DI -- Universit\'e de Moundou 2026
\end{center}

\end{document}
